% Options for packages loaded elsewhere
\PassOptionsToPackage{unicode}{hyperref}
\PassOptionsToPackage{hyphens}{url}
\PassOptionsToPackage{dvipsnames,svgnames,x11names}{xcolor}
%
\documentclass[
  11pt,
  a4paper,
]{ltjsarticle}
\usepackage{amsmath,amssymb}
\usepackage{iftex}
\ifPDFTeX
  \usepackage[T1]{fontenc}
  \usepackage[utf8]{inputenc}
  \usepackage{textcomp} % provide euro and other symbols
\else % if luatex or xetex
  \usepackage{unicode-math} % this also loads fontspec
  \defaultfontfeatures{Scale=MatchLowercase}
  \defaultfontfeatures[\rmfamily]{Ligatures=TeX,Scale=1}
\fi
\usepackage{lmodern}
\ifPDFTeX\else
  % xetex/luatex font selection
\fi
% Use upquote if available, for straight quotes in verbatim environments
\IfFileExists{upquote.sty}{\usepackage{upquote}}{}
\IfFileExists{microtype.sty}{% use microtype if available
  \usepackage[]{microtype}
  \UseMicrotypeSet[protrusion]{basicmath} % disable protrusion for tt fonts
}{}
\makeatletter
\@ifundefined{KOMAClassName}{% if non-KOMA class
  \IfFileExists{parskip.sty}{%
    \usepackage{parskip}
  }{% else
    \setlength{\parindent}{0pt}
    \setlength{\parskip}{6pt plus 2pt minus 1pt}}
}{% if KOMA class
  \KOMAoptions{parskip=half}}
\makeatother
\usepackage{xcolor}
\usepackage[margin=24mm]{geometry}
\setlength{\emergencystretch}{3em} % prevent overfull lines
\providecommand{\tightlist}{%
  \setlength{\itemsep}{0pt}\setlength{\parskip}{0pt}}
\setcounter{secnumdepth}{5}
\ifLuaTeX
  \usepackage{selnolig}  % disable illegal ligatures
\fi
\IfFileExists{bookmark.sty}{\usepackage{bookmark}}{\usepackage{hyperref}}
\IfFileExists{xurl.sty}{\usepackage{xurl}}{} % add URL line breaks if available
\urlstyle{same}
\hypersetup{
  colorlinks=true,
  linkcolor={blue},
  filecolor={Maroon},
  citecolor={Blue},
  urlcolor={blue},
  pdfcreator={LaTeX via pandoc}}

\author{}
\date{}

\begin{document}

\hypertarget{ux89e3ux7b54-06-ux56faux6709ux5024}{%
\section{解答 06 --- 固有値}\label{ux89e3ux7b54-06-ux56faux6709ux5024}}

\hypertarget{section}{%
\subsection{1}\label{section}}

特性方程式は

\[
(2-\lambda)^2-1=0
\]

なので \texttt{λ=3,1}。対応する正規化固有ベクトルは

\[
v_1=\frac1{\sqrt2}(1,1)^T,\quad v_2=\frac1{\sqrt2}(1,-1)^T.
\]

\hypertarget{section-1}{%
\subsection{2}\label{section-1}}

\texttt{Q={[}v\_1\ v\_2{]}}, \texttt{Λ=diag(3,1)} とすれば
\texttt{A=QΛQ\^{}T} です。

\hypertarget{section-2}{%
\subsection{3}\label{section-2}}

\[
A^kx_0=c_13^kv_1+c_2v_2.
\]

\texttt{c\_1≠0} なら大きな \texttt{k} では最大絶対固有値に対応する
\texttt{v\_1} 方向が支配します。

\hypertarget{section-3}{%
\subsection{4}\label{section-3}}

\texttt{Av=λv}, \texttt{Aw=μw} とします。対称性より

\[
v^TAw=(Av)^Tw.
\]

左辺は \texttt{μv\^{}Tw}、右辺は \texttt{λv\^{}Tw}。\texttt{λ≠μ} なので
\texttt{v\^{}Tw=0} です。

\hypertarget{section-4}{%
\subsection{5}\label{section-4}}

任意の \texttt{x} に対し

\[
x^TA^TAx=\|Ax\|^2\ge0.
\]

よって \texttt{A\^{}TA}
は半正定値で、すべての固有値は非負です。この平方根が特異値になります。

\end{document}
